\chapter{Kết luận và Định hướng nghiên cứu}

\section{Tóm tắt các phát hiện và đóng góp chính}
\textbf{Counterfactual Explanation Tree (CET)} được nghiên cứu giới thiệu là một cây quyết định có khả năng gán hành động phản thực hiệu quả cho từng đầu vào trên toàn bộ không gian dữ liệu. Các đóng góp chính bao gồm:
\begin{itemize}
    \item Giới thiệu \textbf{CET (Counterfactual Explanation Tree)}, một cây quyết định có khả năng gán hành động thay đổi kết quả đầu ra một cách hiệu quả cho từng đầu vào trên toàn bộ không gian dữ liệu. CET tận dụng cấu trúc cây quyết định (Decision Tree) để (1) cung cấp quy trình \textbf{minh bạch} (Transparency) trong việc gán hành động, và (2) gán duy nhất một cặp \textit{hành động và lý do} cho mỗi cá thể để đảm bảo tính nhất quán trên toàn bộ tập dữ liệu.
    
    \item Xây dựng phương pháp học CET từ một tập dữ liệu cho trước, giải quyết bài toán tối ưu hóa bằng thuật toán \textbf{Stochastic Local Search}, kết hợp chiến lược pruning dựa trên ràng buộc mục tiêu.
    
    \item Thử nghiệm trên các tập dữ liệu công khai để đánh giá hiệu quả của CET so với các phương pháp hiện có. Ngoài ra, qua các nghiên cứu về người dùng, kết quả cho thấy CET dễ hiểu và trực quan đối với người dùng hơn.
\end{itemize}


\section{Ưu điểm và hạn chế của phương pháp}

\textbf{Ưu điểm:}
\begin{itemize}
    \item \textbf{Minh bạch và nhất quán:} CET tận dụng cấu trúc cây quyết định để gán duy nhất một cặp \textit{hành động + lý do} cho mỗi cá thể, tránh tình trạng gán nhiều hành động hoặc thiếu bao phủ như trong AReS.
    \item \textbf{Phủ toàn cục:} Khác với CE cục bộ, CET đảm bảo phân vùng toàn bộ không gian đầu vào và đưa ra một bản tóm tắt toàn cục đáng tin cậy.
    \item \textbf{Dễ hiểu với người dùng:} Kết quả khảo sát người dùng cho thấy CET trực quan và dễ diễn giải hơn so với các phương pháp tổng hợp toàn cục khác.
\end{itemize}

\textbf{Hạn chế:}
\begin{itemize}
    \item \textbf{Giới hạn về khả năng biểu đạt:} Do chỉ sử dụng một cây quyết định, CET có thể gặp khó khăn trong việc mô tả các quan hệ phi tuyến phức tạp, nơi mà các mô hình ensemble thường vượt trội hơn.
    \item \textbf{Chi phí tính toán:} Việc học CET được mô hình hóa như một bài toán tối ưu hóa toàn cục và giải bằng \textit{Stochastic Local Search} kết hợp với pruning. Do đó, khi dữ liệu có số chiều lớn hoặc nhiều mẫu, chi phí tính toán có thể tăng đáng kể.
\end{itemize}


\section{Ý nghĩa đối với lĩnh vực}
CET mở ra một hướng tiếp cận mới trong việc tổng hợp các giải thích phản thực: vừa đảm bảo \textbf{minh bạch} và \textbf{nhất quán}, vừa có thể áp dụng cho toàn bộ không gian đầu vào. Điều này đặc biệt có ý nghĩa trong các lĩnh vực nhạy cảm như \textbf{tài chính}, \textbf{y tế}, và \textbf{pháp lý}, nơi mà người dùng không chỉ cần dự đoán chính xác mà còn phải hiểu rõ \textit{những hành động khả thi để thay đổi kết quả mô hình}. Bằng việc cung cấp cả góc nhìn cục bộ lẫn toàn cục một cách rõ ràng, CET góp phần thu hẹp khoảng cách giữa \textit{hiệu quả mô hình} và \textit{khả năng giải thích}, giúp tăng niềm tin của người dùng vào các hệ thống học máy.


\section{Định hướng nghiên cứu trong tương lai}
Các hướng nghiên cứu trong tương lai có thể bao gồm:
\begin{itemize}
    \item \textbf{Mở rộng CET cho dữ liệu phi tuyến và nhiều đặc trưng}: áp dụng các kỹ thuật phân nhánh phức tạp hơn hoặc kết hợp với học sâu.
    \item \textbf{Tối ưu hóa hiệu năng}: giảm thời gian tính toán khi áp dụng CET trên các tập dữ liệu lớn.
    \item \textbf{Kết hợp với các loại giải thích khác}: tích hợp CE với SHAP hoặc LIME để tăng cường độ chính xác và khả năng minh bạch (SHAP/LIME có thể giúp xác định đặc trưng quan trọng nhất cần thay đổi, từ đó giảm chi phí tìm kiếm hành động CE).
\end{itemize}
